\subsection{Обоснование необходимости разработки программы}
\label{subsec:development-rationale}

Многопоточная программа использует вычислительные ресурсы через несколько уровней организации
  исполнения. Алгоритм разделяется на задачи, задачи выполняются программными потоками, а потоки ОС
  получают процессорное время. Поэтому одинаковый объём вычислений может занимать разное время
  в зависимости от размера задач, способа их распределения и частоты обращения к общим данным.
  В рамках архитектуры вычислительных систем представляет интерес связь этих программных решений
  с загрузкой ядер, использованием памяти и расходами на планирование.

Для исследования этой связи разрабатывается программа, исполняющая одинаковуе типы задач
  при разных способах организации горутин на потоки ОС. Она позволяет изменять условия по отдельности:
  число горутин, доступный параллелизм, размер пакета данных, способ накопления результата, наличие
  ожидания и использование ресурсов системы. Такое управление необходимо, чтобы различие между вариантами можно было связать
  с конкретным изменением организации исполнения.

Реальный сервис объединяет различные механизмы организации и стуктуризации нагрузки: собственные очереди, библиотеки,
  io-bound и cpu-bound задачи и фоновые операции. Отдельная программа даёт возможность сохранить
  общий алгоритм и одинаковые входные данные в двух режимах организации, а также проверить результат по
  последовательному варианту. При этом измеряется поведение выбранных реализаций на конкретных типах задач, а не универсальное
  преимущество одной модели для любых типов задач.

Практическая область исследования — обработка журналов событий, телеметрии, потока задач. Такие данные
  разбираются, группируются, преобразуются и подготавливаются к передаче или хранению.
  OpenTelemetry Collector описывает связанные с этим ограничения: накопление данных в очередях,
  расход памяти и недостаточную скорость получателя. Увеличение числа параллельных обработчиков
  помогает только тогда, когда оно устраняет действительное ограничение
  производительности~\cite{otel_collector_scaling}.

Назначение разработки состоит в получении воспроизводимых измерений, по которым можно оценить
  полезность увеличения параллелизма и стоимость выбранного способа организации потоков на ядрах ОС.
  Результат представляет интерес и при выборе размера пакетов и числа одновременно выполняемых
  задач в программах обработки данных.

\subsection{Модели многопоточности и ограничения параллельного ускорения}
\label{subsec:threading-theory}

Конкурентное выполнение означает, что несколько задач находятся в процессе обработки и могут
  выполняться независимо. Параллельное выполнение означает одновременное исполнение на нескольких
  вычислительных ресурсах. На одном аппаратном потоке задачи могут чередоваться, тем самым оставаясь
  конкурентными. Поэтому число созданных программных потоков само по себе не характеризует
  достигнутый параллелизм.

В модели 1:1 каждому пользовательскому потоку соответствует отдельный поток ОС. Его исполнение
  управляется планировщиком операционной системой наряду с другими системными потоками. Такая организация
  связывает увеличение числа пользовательских потоков с ресурсами, необходимыми для системных
  потоков: их состоянием, стеками и операциями планирования.

В модели M:N множество пользовательских потоков распределяется средой выполнения между предоставленным набором
  потоков ОС. Пользовательский поток, ожидающий событие, например: блокировка примитивами синхронизации,
  ожидание ресурсов, вызов ввода-вывода, может уступить возможность исполнения на потоке
  другой готовой задаче без создания для неё нового системного потока, тем самым снимая часть накладных расходов. Однако это требует
  дополнительного планировщика и его согласованность с работой планировщика ОС. В Go такое распределение реализуется
  средой выполнения, описанной в подразделе~\ref{subsec:go-runtime}~\cite{go_runtime_hacking}.

Освобождение потока для другой горутины возможно при ожиданиях, обслуживаемых средой выполнения.
  Блокирующий системный вызов может занять поток ОС; продолжение других горутин в этом случае
  обеспечивается другими потоками среды~\cite{go_runtime_hacking}.

Ни одна из моделей не устраняет зависимости между данными, ограничения памяти и необходимость
  синхронизации. Гонка данных (\texttt{data race}) возникает при неупорядоченных средствами
  синхронизации обращениях разных горутин к одной области памяти, если хотя бы одно обращение
  является записью и обращения не защищены атомарными операциями~\cite{go_memory_model}.
  Состояние гонки (\texttt{race condition}) — ошибка, при которой корректность результата зависит
  от взаимного порядка или момента выполнения конкурентных операций. Такая ошибка возможна
  и при отсутствии гонки данных. Общий мьютекс обеспечивает взаимное исключение при обновлении
  статистики, но одновременно ограничивает параллелизм этой части обработки.
  Если несколько ядер одновременно читают большой набор данных, ограничением может стать уже
  пропускная способность памяти. Переключение модели не увеличивает число ядер и не отменяет необходимость синхронизации.

Ограничение ускорения из-за последовательной части известно как закон Амдала~\cite{amdahl1967}.
  Для фиксированного объёма работы его распространённая запись имеет вид
\begin{equation}
  S_{\mathrm{ideal}}(p) = \frac{1}{s + \frac{1-s}{p}},
  \label{eq:amdahl-speedup}
\end{equation}
где $s$ — доля последовательной работы во времени исходного выполнения, а $p$ — число равноценных
  параллельных вычислительных ресурсов. При $s > 0$ увеличение $p$ ограничивает ускорение значением
  $1/s$. Формула~\eqref{eq:amdahl-speedup} предполагает отсутствие дополнительных расходов.
  Реальная программа также затрачивает время на очереди, синхронизацию, планирование и объединение
  результатов, поэтому формула служит ориентиром для объяснения предела масштабирования.

На неоднородных процессорах ресурсы неравноценны: производительные и энергоэффективные ядра
  различаются возможностями, а аппаратные потоки одного ядра используют общие ресурсы.
  Поэтому ни число горутин, ни значение \texttt{GOMAXPROCS} нельзя непосредственно подставлять
  в формулу как точное число одинаковых независимых процессоров.

% TODO: вставить схему отображения пользовательских потоков на потоки ОС для моделей 1:1 и M:N.
% Показать ядра отдельным уровнем; не смешивать потоки ОС с аппаратными потоками процессора.

\subsection{Прикладные исследования и выбор задач}
\label{subsec:workload-rationale}

Выбор задачи определяет, какие свойства проявятся в измерениях. Только длительные
  независимые вычисления позволяют оценить использование CPU, но могут скрыть относительную
  стоимость планирования потоков. Только ожидание событий, напротив, мало характеризует выполнение
  параллельных вычислений. Поэтому исследование охватывает различные типы нагрузок в обработке данных.

В работе Shinjuku исследуется планирование запросов с различным временем обслуживания, включая
  точечные и диапазонные запросы RocksDB. Важная для данного проекта проблема — ожидание короткого
  запроса за длинным~\cite{shinjuku2019}. В Concord рассматривается соотношение стоимости
  планирования, пропускной способности и задержек; оценка включает LevelDB~\cite{concord2023}.
  Эти исследования обосновывают интерес к неоднородным задачам и задержкам отдельных операций,
  но сравнивают другие системы планирования в рамках баз данных.

Набор PARSEC включает прикладные многопоточные программы с разными рабочими наборами, локальностью
  данных и способами синхронизации~\cite{parsec2008}. Такой подход показывает значение разнообразия
  нагрузки для архитектурного исследования. В курсовом проекте используется небольшой управляемый
  набор сценариев, позволяющий объяснить результат каждой серии на доступных системах.

Первый сценарий — разбор записей журналов и накопление статистики. Запись содержит идентификатор
  сервиса, код ответа, длительность операции и размер сообщения. Обработка пакета включает разбор
  полей, фильтрацию и подсчёт сумм и гистограммы. Изменение числа записей в пакете меняет соотношение
  полезной работы и расходов на выдачу задач. Данные размещаются в памяти до начала измерения,
  чтобы основная серия не зависела от скорости чтения накопителя.

Для этой нагрузки рассматриваются два способа накопления результата: обновление общей статистики
  используя средства синхронизации и локальный подсчёт с последующим объединением. Оба способа дают один результат,
  но различаются частотой необходимой синхронизации. Каждый из них сравнивается в обеих моделях потоков,
  что позволяет исследовать синхронизацию как отдельный фактор.

Второй сценарий — сжатие независимых пакетов данных. Он связан с подготовкой журналов к хранению
  или передаче и даёт более длительную обработку отдельного пакета. Практическим примером
  параллельного сжатия служит pigz, распределяющий сжатие блоков между потоками~\cite{pigz_manual}.
  В данном проекте используются независимые блоки и общий алгоритм сжатия; полный конвейер pigz
  не воспроизводится. Размер блока, уровень сжатия и свойства содержимого влияют на стоимость
  обработки, поэтому сохраняются одинаковыми внутри сравниваемой пары.

Смесь пакетов разных размеров позволяет исследовать лёгкие и тяжёлые задачи. Класс задачи
  определяется объёмом реальной обработки. Доля тяжёлых задач по их количеству отличается
  от их доли в байтах и затратах CPU; эти характеристики рассматриваются отдельно.
  Поиск по ключу и обработка диапазона представляют альтернативный способ формирования
  неоднородной нагрузки, однако добавляют зависимость от устройства индекса и локальности памяти.

Для воспроизводимости содержимое данных и порядок задач задаются отдельными начальными состояниями
  псевдослучайного генератора: \texttt{data\_seed} и \texttt{order\_seed}. Сначала формируется
  точное число задач каждого класса, затем готовый список перемешивается. Это сохраняет состав
  нагрузки при изменении порядка. Разные значения \texttt{order\_seed} позволяют оценить влияние
  расположения тяжёлых задач, а повторение одного списка — изменчивость времени запуска.
  Порядок исполнения остаётся зависимым от планировщиков даже при одинаковом порядке выдачи.

Дополнительный сценарий чередует обработку пакета с ожиданием таймера. Он позволяет наблюдать
  расход потоков и памяти при большом числе ожидающих горутин. Заданная задержка моделирует
  ожидание, но не характеризует реальную сеть или накопитель. Поэтому результаты этой серии
  рассматриваются отдельно от результатов вычислительной нагрузки.

% TODO: при окончательном выборе нагрузки согласовать сценарии с требованиями раздела 4 и матрицей
% раздела 5. Поиск остаётся альтернативой сжатию; реальный TCP-обмен — возможным расширением.

\subsection{Технологии программирования для решения поставленных задач}
\label{subsec:implementation-technologies}

Для реализации выбран Go, поскольку он предоставляет горутины, средства синхронизации и управление
  параметрами среды выполнения в одной платформе. Это позволяет использовать общую
  полезную функцию в сравниваемых режимах и уменьшить число различающихся компонентов программы.
  Выбор языка здесь подчинён организации эксперимента; сравнение языков в задачу работы не входит.

Последовательный режим \texttt{seq} выполняет те же операции без параллельной очереди и будет служить
  эталоном результата и времени для сравнения. В режиме \texttt{mn} используются обычные горутины, распределяемые
  средой Go между потоками ОС. В режиме \texttt{pinned} каждая долгоживущая горутина обработки
  закрепляется за отдельным системным потоком вызовом \texttt{runtime.LockOSThread}.
  Закрепление снимается парным вызовом \texttt{runtime.UnlockOSThread}~\cite{go_runtime_package}.

Режим \texttt{pinned} приближает отображение 1:1, сохраняя планировщик Go. Для исполнения
  пользовательского кода по-прежнему нужен ресурс P, поэтому число закреплённых горутин может
  превышать число одновременно выполняющихся горутин. Режимы \texttt{mn} и \texttt{pinned}
  используют одинаковое число горутин обработки и одну организацию очереди. Сравнение с
  \texttt{seq} показывает эффект параллелизации в целом, а сравнение двух параллельных режимов
  характеризует влияние закрепления при общей программной основе~\cite{go_runtime_hacking}.

При общей политике статистика сценария защищается мьютексом. При локальной политике каждый вызов
  \texttt{Run} формирует отдельный частичный результат, который передаётся модулю эксперимента
  для объединения. Исполнитель не изменяет правила накопления статистики при переключении режима.
  Средства передачи задач обеспечивают согласованную работу очереди.
  Для сжатия подходит пакет стандартной библиотеки \texttt{compress/flate}, реализующий DEFLATE
  и позволяющий задавать уровень сжатия~\cite{go_flate_package}. Одинаковая библиотека в обоих
  режимах сохраняет алгоритм преобразования данных.

Профилирование CPU, памяти и блокировок, а также трассировка исполнения дополняют измерение времени.
  Они помогут различить полезную обработку, ожидание и работу среды выполнения.
  Диагностические запуски проводятся отдельно, поскольку сбор подробных профилей сам создаёт
  дополнительные расходы~\cite{go_diagnostics}. Их назначение — проверить объяснение наблюдаемой
  разницы, а не заменить основную серию измерений.

\subsection{Связь архитектуры системы с разрабатываемым программным обеспечением}
\label{subsec:architecture-software-link}

При обработке независимых пакетов программа может использовать несколько ядер одновременно.
  Однако число задач, горутин и потоков ОС описывает разные уровни исполнения. Одна горутина
  последовательно обрабатывает несколько задач, а число потоков среды выполнения включает
  также служебные потоки. Поэтому оценка использования аппаратуры требует учитывать эти величины
  раздельно и сопоставлять их с доступными логическими CPU.

Мелкие пакеты увеличивают частоту операций с очередью относительно объёма обработки. Крупные
  пакеты изменяют рабочий набор и продолжительность занятости горутины одной задачей.
  Это связывает программное разбиение данных с кэшем, основной памятью и возможностями балансировки.
  Большой пакет может уменьшить расходы выдачи задач, но одновременно увеличить задержку
  последующих задач. Поэтому один размер пакета не характеризует все условия выполнения.

При общей статистике ядра обращаются к одним изменяемым данным. Локальное накопление уменьшает
  частоту совместных обновлений, но требует дополнительной памяти и итогового объединения.
  Выбор способа накопления, таким образом, изменяет как взаимодействие потоков, так и характер
  обращений к памяти. Время объединения является частью получения полезного результата.

На выбранных процессорах размещение потоков по P- и E-ядрам выполняет ОС. Закрепление горутины
  через \texttt{LockOSThread} фиксирует соответствие системному потоку, но не его положение на CPU.
  Значение \texttt{GOMAXPROCS} также не задаёт класс используемых ядер. Эти особенности ограничивают
  интерпретацию зависимости времени от числа горутин на неоднородных системах.

Основная пара измерений выполняется при одинаковой конфигурации одной машины. Повторение на второй
  платформе позволяет проверить сохранение характера зависимостей. Поскольку одновременно меняются
  процессор, память, ОС и условия охлаждения, межплатформенная разница характеризует системы целиком.
  Она не является отдельной оценкой влияния x86, ARM, Linux или macOS.

\subsection{Проверяемые гипотезы и показатели исследования}
\label{subsec:research-hypotheses}

Теоретические положения задают следующие гипотезы для экспериментальной проверки:
\begin{enumerate}
  \item при уменьшении пакета возрастает относительное влияние очереди и планирования на время
    обработки фиксированного объёма данных;
  \item увеличение числа горутин сверх доступного параллелизма может повышать расходы,
    причём закрепление способно увеличивать число системных потоков и потребление ресурсов;
  \item общая блокировка ограничивает масштабирование сильнее, чем локальное накопление,
    и может скрывать различие способов отображения потоков;
  \item смесь лёгких и тяжёлых задач и их порядок влияют на задержки задач;
  \item при большом числе ожидающих задач M:N может обеспечивать сопоставимую обработку
    с меньшим числом потоков ОС, чем вариант с закреплением.
\end{enumerate}

Гипотезы допускают отсутствие эффекта или изменение его направления в зависимости от параметров.
  Их проверка опирается на повторные измерения и дополнительные показатели; отдельный медленный
  запуск или само число созданных горутин не подтверждают причину изменения производительности.

Основной показатель — время $T$ обработки заданного набора данных, включая выдачу задач
  и получение объединённого результата. Для нагрузки объёмом $V$ пропускная способность равна
\begin{equation}
  Q = \frac{V}{T},
  \label{eq:throughput}
\end{equation}
где $V$ измеряется в записях либо байтах в зависимости от сценария. Единица $Q$ соответственно
  задаётся в записях в секунду или байтах в секунду. Сравнение корректно при одинаковом содержании
  полезной операции и одинаковом объёме внутри каждой пары запусков.

Ускорение относительно последовательного режима определяется как
\begin{equation}
  S = \frac{T_{\mathrm{seq}}}{T},
  \label{eq:measured-speedup}
\end{equation}
где $T_{\mathrm{seq}}$ — время последовательной обработки того же набора на той же машине.
  Этот показатель включает эффект выбранной параллельной организации целиком.

Для непосредственного сопоставления двух параллельных режимов используется отношение
\begin{equation}
  R = \frac{T_{\mathrm{pinned}}}{T_{\mathrm{mn}}}.
  \label{eq:threading-ratio}
\end{equation}
При $R > 1$ обычные горутины выполняют заданную работу быстрее закреплённых, при $R < 1$ — медленнее.
  Значение около единицы оценивается с учётом разброса повторных запусков и само по себе
  не доказывает равенства реализаций.

Дополнительные показатели — число живых потоков ОС, память процесса и задержки отдельных задач.
  Память процесса не равна размеру кучи Go, а число потоков ОС не равно числу горутин.
  В отдельной серии задержка задачи отсчитывается от момента начала её передачи исполнителю
  до завершения обработки. Она включает ожидание свободного места в ограниченной очереди,
  нахождение в очереди и выполнение задачи. Время пребывания в очереди и длительность самой
  обработки при необходимости рассматриваются как отдельные показатели.
  Для неоднородной нагрузки задержки анализируются отдельно по классам задач, чтобы длительные
  операции не скрывали поведение коротких. Порядок измерений и обработки результатов раскрывается
  в разделе~\ref{sec:results}.
